\documentclass[a4paper, 12pt]{article}

\usepackage[utf8]{inputenc}
\usepackage[portuguese,brazil]{babel}
\usepackage{amsmath, amssymb, amsfonts}
\usepackage[a4paper,left=3cm,right=2cm,top=3cm,bottom=2cm]{geometry}
% Caso não queira o 'geometry', pode-se usar:
% \usepackage{a4wide}
\usepackage{tabularx}
\usepackage{graphicx}
\usepackage{pslatex}
\usepackage{indentfirst}
\usepackage{setspace}
\usepackage{verbatim}

\begin{document}

\thispagestyle{empty}

\begin{minipage}{0.15\linewidth}
\begin{flushleft}
\includegraphics[width=\linewidth]{UFRN}
\end{flushleft}
\end{minipage}\hfill
\begin{minipage}{0.65\linewidth}
\centering Universidade Federal do Rio Grande do Norte\\
Centro de Tecnologia\\
Coordenação de Engenharia de Computação
\end{minipage}\hfill
\begin{minipage}{0.15\linewidth}
\begin{flushright}
\includegraphics[width=0.6\linewidth]{EngComp}
\end{flushright}
\end{minipage}


%No caso de alunos de engenharia mecatrônica, como não há logotipo do curso, então pode-se adotar o cabeçalho inicial como:
%\begin{center}
%Universidade Federal do Rio Grande do Norte\\
%Centro de Tecnologia\\
%Coordenação de Engenharia Mecatrônica
%\end{center}
% Apague o trecho acima, que possui os logotipos da UFRN e da coordenação de eng. computação, substituindo por esse que está
%comentado!

\vfill

\begin{center}
\LARGE{\bf Relatório de Estágio Supervisionado}

\Large{Nome da Empresa}
\end{center}

\normalsize
\vfill

\begin{flushleft}
{\bf Aluno(a):} nome do aluno \\
{\bf Supervisor:} nome do supervisor\\
{\bf Orientador acadêmico:} nome do orientador 
\end{flushleft}

\vfill

\begin{center}
Natal, RN, \today.
\end{center}

\newpage

\pagenumbering{roman}

\begin{center}
{\bf \Large Agradecimentos}
\end{center}

São menções que o autor faz a pessoas e/ou instituições das quais eventualmente recebeu apoio e que concorreram de maneira relevante para o desenvolvimento do relatório. Lembrar que este item é opcional.

\newpage

\tableofcontents

\newpage

% Colocar somente se necessáro...
\listoffigures
\addcontentsline{fachada_sinfo}{section}{Lista de Figuras}
\addcontentsline{anexo_sinfo}{section}{Lista de Figuras}


\newpage

% Colocar somente se necessáro...
\listoftables
\addcontentsline{toc}{section}{Lista de Tabelas}


\newpage

\section*{Resumo}
\addcontentsline{toc}{section}{Resumo}

É a apresentação concisa do texto, destacando seus aspectos de maior relevância. Na elaboração do resumo deve-se:

\begin{itemize}
\item apresentar o resumo precedendo o texto, em entrelinhamento menor e em página distinta;
\item redigir em um único parágrafo;
\item redigir com frases completas e não com sequência de títulos;
\item expressar na primeira frase do resumo o assunto tratado, situando-o no tempo e no espaço, caso o título do trabalho não seja suficientemente explícito;
\item dar preferência ao uso da terceira pessoa do singular;
\item evitar o uso de citações bibliográficas;
\item ressaltar os objetivos, os métodos, os resultados e as conclusões do trabalho;
\item elaborar o resumo de relatórios com, no máximo, 500 palavras.
\end{itemize}


\newpage
\pagenumbering{arabic}
\onehalfspacing

\section{Ambiente de Estágio}

Pode-se usar outro título para esta seção ou iniciar com outro assunto pertinente. Em geral, é necessário que o relatório de estágio aborde o ambiente onde o aluno atuou. Nisto incluem-se uma descrição do histórico da empresa, a sua descrição física (localização, entre outros) e quaisquer outras informações pertinentes. Lembrar que algumas informações podem não ser de domínio público. Nestas situações, o supervisor do ambiente de trabalho poderá direcionar melhor a redação.

% Caso deseje...
% \newpage

\section{Atividades Desenvolvidas}

Descrição detalhada das atividades realizadas pelo estagiário, incluindo metodologia, cronogramas adotados, referências pesquisadas. Normalmente esta será a maior seção do relatório. Para uma melhor organização da seção, inclua subseções com 
\verb+\subsection+.

% Caso deseje...
% \newpage

\section{Outras Informações}

Caso queira descrever alguma particularidade técnica ou outra informação que não caiba nas atividades desenvolvidas, você pode criar seções específicas para isso.

Trantando especificamente deste \emph{template} em {\LaTeX}, você pode editá-lo da maneira que desejar. Apenas mantenha a estrutura original para que os relatórios de estágio sigam um padrão. Se você usa o ambiente linux, então a geração do arquivo PDF se dá através de:

\begin{verbatim}
latex nome_arquivo.tex
latex nome_arquivo.tex
dvips nome_arquivo.dvi
ps2pdf nome_arquivo.ps
\end{verbatim}
ou alternativamente por:
\begin{verbatim}
pdflatex nome_arquivo.tex
\end{verbatim}

No final você terá um arquivo PDF chamado \verb+nome_arquivo.pdf+.

Se você usa ambiente Windows$\copyright$, então use o \emph{MikTex}, disponível em:
\begin{center}
\texttt{http://miktex.org/}
\end{center}

A grande maioria das questões técnicas sobre o {\LaTeX} pode ser obtida pela internet ou em livros, como o de Kopka e Daly 
\cite{livtex}. Na Biblioteca Central Zila Mamede (UFRN) há alguns exemplares de outros livros que tratam deste assunto.

Caso alguem prefira usar outro editor de texto, em breve providenciaremos \emph{templates} para arquivos DOCX e ODT.

Os únicos exemplos apresentados aqui serão os da inserção de figura e também de tabela. Em cada uma delas, o {\LaTeX} necessita de um ambiente específico. Veja o exemplo para figura a seguir:

\begin{verbatim}
\begin{figure}[htb!]
\begin{center}
\includegraphics[width=5cm]{tux}
\end{center}
\caption{Tux no DCA!}
\label{tux}
\end{figure}
\end{verbatim}

que reproduz a figura~\ref{fachada_sinfo}.
\begin{figure}[htb!]
\begin{center}
\includegraphics[width=5cm]{fachada_sinfo}
\end{center}
\caption{Tux no DCA!}
\label{tux}
\end{figure}

que reproduz a figura~\ref{anexo_sinfo}.
\begin{figure}[htb!]
\begin{center}
\includegraphics[width=5cm]{anexo_sinfo}
\end{center}
\caption{Tux no DCA!}
\label{tux}
\end{figure}

Já a tabela é obtida da seguinte forma:
\small
\begin{verbatim}
\begin{center}
\begin{tabular}{|r|l||c|rrr|c|c|} \hline
Position & Club         & Games & W  & T  & L  & Goals & Points\\ \hline\hline
1        & Amesville Rockets & 33    & 19 & 13 & 1  & 66:31 & 51:15 \\ \hline
2        & Ralston Regulars  & 33    & 3  & 11 & 19 & 37:74 & 17:49 \\ \hline
\end{tabular}
\end{center}
\end{verbatim}
\normalsize

\begin{center}
\begin{tabular}{|r|l||c|rrr|c|c|} \hline
Position & Club              & Games & W  & T  & L  & Goals & Points\\ \hline\hline
1        & Amesville Rockets & 33    & 19 & 13 & 1  & 66:31 & 51:15 \\ \hline
2        & Ralston Regulars  & 33    & 3  & 11 & 19 & 37:74 & 17:49 \\ \hline
\end{tabular}
\end{center}

% Caso deseje...
\newpage

\section{Conclusões}

Suas impressões sobre o estágio realizado. Faça menção à atividade de estágio, caso esteja matriculado nela. É no final desta folha que há um espaço para sua assinatura.

Última dica: ao final, revise todo o documento PDF gerado pelo {\LaTeX} pois, em algumas situações, ele não promove a separação silábica correta em português.

\vfill

\begin{center}
\begin{tabular}{c}
\rule{10cm}{1pt} \\ 
Nome do estagiário
\end{tabular}
\end{center}

\newpage

\section*{Comentários do Supervisor}
\addcontentsline{toc}{section}{Comentários do Supervisor do Estágio na Empresa}

Espaço destinado para os comentários do supervisor do estagiário na empresa. Pode-se deixar a folha em branco para que ele faça anotações à mão livre. Há um espaço abaixo para a assinatura dele.

\vfill

\begin{center}
\begin{tabular}{c}
\rule{10cm}{1pt} \\ 
Nome do supervisor
\end{tabular}
\end{center}

\newpage

\section*{Comentários do Orientador Acadêmico}
\addcontentsline{toc}{section}{Comentários do Orientador do Estágio na UFRN}

Espaço destinado para os comentários do orientador do estagiário na UFRN. Também pode-se deixar a folha em branco para que ele faça anotações à mão livre.  Há um espaço abaixo para a assinatura dele.

\vfill

\begin{center}
\begin{tabular}{c}
\rule{10cm}{1pt} \\ 
Nome do orientador
\end{tabular}
\end{center}

\newpage

\singlespacing

\addcontentsline{toc}{section}{Referências}
\begin{thebibliography}{99}

\bibitem{livtex} Helmut Kopka; Patrick W. Daly. {\em A Guide to \LaTeX\ and Electronic Publishing}. 4$^\mathrm{a}$ edição. Editora Addison-Wesley. 2004.
\end{thebibliography}


\end{document}


